\chapter{Фрагменты исходного кода параллельного расчётного модуля}
\label{app:code}

В листинге~\ref{lst:threads} показана установка числа потоков по параметру командной строки, в листинге~\ref{lst:main} --- основной цикл по времени со слиянием пространственных циклов, а в листинге~\ref{lst:mean} --- вычисление средней температуры с редукцией. Полный исходный код находится в файле \texttt{heat3d.cpp} репозитория программного комплекса.

\begin{lstlisting}[language=C++, caption={Установка числа потоков OpenMP}, label=lst:threads]
int nThreads = 1;
#ifdef _OPENMP
    if (p.threads > 0) omp_set_num_threads(p.threads);
    nThreads = omp_get_max_threads();
#endif
\end{lstlisting}

\begin{lstlisting}[language=C++, caption={Основной цикл по времени}, label=lst:main]
const int NIJ = (p.nx - 2) * (p.ny - 2);
for (int step = 0; step <= NSTEPS; ++step) {
    if (step % p.save_every == 0) {
        /* write T_center, T_mean and slices to CSV */
    }
    if (step == NSTEPS) break;

    #pragma omp parallel for schedule(static)
    for (int ij = 0; ij < NIJ; ++ij) {
        const int i = 1 + ij / (p.ny - 2);
        const int j = 1 + ij % (p.ny - 2);
        for (int k = 1; k < p.nz - 1; ++k) {
            double Tijk = T[idx(i, j, k, p.ny, p.nz)];
            double d2x = T[idx(i+1, j, k, p.ny, p.nz)] - 2*Tijk
                       + T[idx(i-1, j, k, p.ny, p.nz)];
            double d2y = T[idx(i, j+1, k, p.ny, p.nz)] - 2*Tijk
                       + T[idx(i, j-1, k, p.ny, p.nz)];
            double d2z = T[idx(i, j, k+1, p.ny, p.nz)] - 2*Tijk
                       + T[idx(i, j, k-1, p.ny, p.nz)];
            T_new[idx(i, j, k, p.ny, p.nz)] =
                Tijk + rx*d2x + ry*d2y + rz*d2z;
        }
    }
    applyBC(T_new);      // parallel loops over 3 pairs of faces
    std::swap(T, T_new);
    time += dt;
}
\end{lstlisting}

\begin{lstlisting}[language=C++, caption={Средняя температура по объёму}, label=lst:mean]
double meanTemp(const Grid3D& T) {
    double sum = 0.0;
    const int n = (int)T.size();
    #pragma omp parallel for reduction(+:sum) schedule(static)
    for (int q = 0; q < n; ++q) sum += T[q];
    return sum / (double)T.size();
}
\end{lstlisting}

\chapter{Сборка программного комплекса и запуск вычислительных экспериментов}
\label{app:build}

В листинге~\ref{lst:build} приведены команды сборки расчётного модуля для Linux и Windows, а в листинге~\ref{lst:run} --- команды запуска расчёта с заданным числом потоков, вычислительных экспериментов и автоматических тестов.

\begin{lstlisting}[language=bash, caption={Сборка расчётного модуля}, label=lst:build]
# Linux (GCC)
g++ -O2 -fopenmp -std=c++17 heat3d.cpp -o heat3d
# Windows (MinGW-w64, static linking)
x86_64-w64-mingw32-g++-posix -O2 -fopenmp -std=c++17 \
    -static heat3d.cpp -o heat3d.exe
# Linux: engine + GUI as a single executable (PyInstaller)
./build_linux.sh
\end{lstlisting}

\begin{lstlisting}[language=bash, caption={Запуск расчёта, экспериментов и тестов}, label=lst:run]
# simulation on a 100^3 grid with 4 threads
./heat3d --nx 100 --ny 100 --nz 100 --t_end 20000 --threads 4
# all experiments and figures of chapter 3
python nir_7sem/experiments.py
# speedup measurement for several grids
python report_plots.py --bench --threads 1 2 3 4 --nx 50 100 150
# automated tests of the parallel version
python -m pytest tests/integration/test_openmp.py -v
\end{lstlisting}
